\section{Decision Accuracy across the Complete Budget Catalogue}\label{sec:r18catalogue}
A final-budget comparison alone does not determine which computed procedure
meets an economic tolerance at the lowest cost. We therefore use every
prespecified menu stage. In each economic cell, the catalogue contains the
reference and the five methods at each of three stages. Write $V_i$ for
candidate $i$'s mean payoff over the same complete sixteen-stream population
and 384 queries. The target in this section is
$g_i=\max_j V_j-V_i$ over these sixteen candidates. This is a catalogue of
computed vector decisions, not the full continuous vector-action set or the
adapted diffusion class.

The original simultaneous intervals bound five reference-relative gains and
four direct NBO contrasts at each stage. If $L_{ab}\leq V_a-V_b\leq U_{ab}$,
put directed edges $b\to a$ of length $U_{ab}$ and $a\to b$ of length
$-L_{ab}$. Let $D_{ij}$ be the shortest-path length from $i$ to $j$.
The same confidence event then gives
\begin{equation}\label{eq:r18cataloguegap}
 \max_j\{-D_{ji}\}\ \leq\ g_i\ \leq\ \overline g_i:=\max_j D_{ij}.
\end{equation}
Proposition~\ref{prop:r18closure} proves that these bounds are sharp in the
uncertainty set defined by the supplied intervals. No new simulation,
normal approximation, or additional probability allocation is used.
For a tolerance $\tau$, let $\mathcal C_\tau=\{i:\overline g_i\leq\tau\}$.
Selecting the least recorded cost within $\mathcal C_\tau$ preserves the
simultaneous accuracy guarantee. It identifies the cheapest \emph{certified}
candidate, not necessarily the cheapest candidate whose unknown true regret
is below $\tau$.

\begin{table}[htbp]\centering
\caption{Least Accounted Cost with a Catalogue Accuracy Certificate}\label{tab:r18catalogue}
\small\setlength{\tabcolsep}{4pt}
\begin{tabular}{lrlrrr}\toprule
Design & $d$ & Certified choice & $\overline g$ & Seconds & NBO (3): $\overline g$\\\midrule
Low & 10 & SAA (1) & 0.03935 & 3.47 & 0.03883\\
Low & 50 & SAA (1) & 0.03563 & 3.97 & 0.05945\\
Quarterly & 10 & Raw actor (2) & 0.7543 & 4.85 & 0.1848\\
Quarterly & 50 & Vector (3) & 0.7678 & 11.15 & 0.5235\\
Long & 10 & SAA (3) & 0 & 30.42 & 1.984\\
Long & 50 & SAA (3) & 0.4893 & 67.04 & 2.688\\
Intermediate & 10 & Vector (2) & 0.8382 & 6.52 & 0.4282\\
Intermediate & 50 & NBO (3) & 0.9140 & 12.70 & 0.9140\\
\bottomrule\end{tabular}
\par\vspace{4pt}\begin{minipage}{.97\textwidth}\footnotesize
Notes: Regret bounds are in units of $10^{-4}$; the target is at most one. Parentheses identify the original stage. The catalogue is the reference plus fifteen method-stage procedures, not the full action space. Early-stage seconds are prefix allocations; stage-three seconds are complete construction/query clocks. Shared assessment and the complete comparative bill are reported separately. The selected cost is least among certified candidates only.
\end{minipage}\end{table}


The fifty-dimensional Intermediate design gives a positive NBO accuracy--work
result against the complete budget catalogue. NBO's final-stage regret upper
bound is $0.000091391$ (rounded upward), below $10^{-4}$, at mean construction
and query cost $12.6993$ seconds (rounded upward). It is the least-cost
candidate with this certificate. The cheaper final Raw and vector actors
have regret lower bounds above $0.0001874$ and $0.0001064$, respectively.
The cheaper four-path SAA and second-stage NBO candidates remain unresolved:
their regret intervals straddle $10^{-4}$. Thus the certified choice does
not establish actual minimum work over these unresolved alternatives.

Table~\ref{tab:r18catalogue} retains the choices in every other design.
The final NBO candidate is within $10^{-4}$ of the best catalogue payoff
in six cells. Both Long cells instead select final-stage SAA; elsewhere a
cheaper Raw actor, vector predictor, or low-budget SAA has the certificate.
The lower-cost alternatives are part of the result, not exclusions from
the comparison.

Early-stage costs are the original construction-prefix allocations, not
independent clocks for a prospectively stopped procedure. Final-stage costs
are measured complete three-stage construction and query processes. The
shared confirmation process is charged separately and in full; adding this
same bill to each candidate does not alter its ordering. The total cost of
having conducted the comparison includes \emph{all} five construction
processes and that shared assessment, not just the selected candidate.
Section~\ref{app:r18record} reports these complete bills and every candidate.
The closure and resulting selection are a post-freeze deterministic
reanalysis of an already simultaneous family, not a newly executed stopping
rule or a fresh confirmatory experiment.
